% ===========================================================================
% sections_v2/related_work.tex
% Drop-in replacement for \section{Related Work} (Sec. II) of
%   ~/mycode/dcdtuni/docs/manuscripts/paper2/main.tex
% Written 2026-09-15 against ddl/docs/FINDINGS_2026-09-15.md Sec. 5 and Sec. 7.
% Every citation key below either already exists in references.bib or is added,
% verified, in sections_v2/references_additions.bib.
% Requires: amsmath, amssymb, booktabs, cleveref (all already loaded by main.tex).
% Cross-file labels: this file defines sec:rw-concurrent (used by
% sections_v2/contributions.tex) and uses sec:contributions (defined there).
% If either file is dropped in alone, replace the missing \cref with the
% plain citation named in the neighbouring comment.
% NOTE: Table I is now a full-width float (table*) because it gained two columns.
%       If a single-column float is required, delete the last column and the
%       entries in it; the prose does not depend on it.
% ===========================================================================
\section{Related Work}
\label{sec:related}

The question ``do differentiable simulators give better policy gradients?'' has
already been asked, and the answer is not the contribution of this paper.
Suh et al.\ showed that stiffness and discontinuities in a physical system
compromise the first-order (pathwise) estimator.  They analysed the failure through
bias and variance.  They proposed an $\alpha$-order estimator, $\alpha\in[0,1]$,
that blends exact first-order gradients with zeroth-order ones
\cite{suh2022diffsim}.  Onoda et al.\ re-examined whether that bias is the
primary obstacle and what minimal fix suffices.  DDCG is a lightweight test that
switches estimators in nonsmooth regions under a single hyperparameter.  IVW-H is a
per-step inverse-variance-weighted hybrid that stabilises the variance without
explicit discontinuity detection.  Their conclusion is that careful
variance control, not bias removal, dominates in practical deployments
\cite{onoda2026ddcg}.  Scheffler et al.\ proved that first-order estimators of a
randomized-smoothing surrogate are biased whenever the underlying function is
discontinuous.  They traced the bias to neglected discontinuity contributions and
derived the corrected first-order formula that eliminates it.  They obtained
estimators that reduce variance by several orders of magnitude over zeroth-order
baselines.  These remain unbiased exactly where standard first-order estimators fail
\cite{scheffler2026cfobg,montaut2022collision}.  For hybrid
discrete--continuous action spaces, Alvo et al.\ combine pathwise and
score-function terms in a mixed estimator that \emph{preserves unbiasedness},
and show that its cross term becomes negligible near a discrete best response
\cite{hpo2026}.  On the language-model side, Hu et al.\ fine-tune an LLM by
amortized Bayesian inference so that it samples an intractable posterior over
discrete sequences rather than maximising a scalar reward.  They offer
distribution matching as an alternative to reward-maximising policy optimisation
\cite{hu2023abs}.

This paper therefore contributes no estimator, no bound on the pathwise bias and
no unbiasedness result; all three exist in the work just cited.  What is missing
there, and what we supply, is the object the estimators act on.  In every one of
them the non-smoothness is a property of the \emph{dynamics}.  The sources are
contact, a status switch and stiffness.  The design or action space the model
emits is already continuous, or already discrete with a relaxation whose physical
meaning is not at issue.  Our setting inverts this: the non-smoothness is in the
\emph{design} the language model emits, and the relaxed object must first be a
system a production solver can simulate at all.  That requirement --- physical
closure of the relaxation --- is a modelling object, not an estimator
correction, and it is the whole of our first contribution
(\cref{sec:contributions}).  Our estimator-side results are
consequently empirical and domain-specific.  We measure that the pathwise bias is
\emph{decision-type dependent}.  We report that a fixed-weight convex
mixture of the two signals is dominated by both of its extremes, because the
pathwise term's error is dominated by scale rather than offset.  That is why we
calibrate it affinely instead.  The unbiased mixed estimator of \cite{hpo2026}
and the inverse-variance weighting of \cite{onoda2026ddcg} are the right objects
for the general problem; our measurement says why neither is what our data asks
for.

\subsection{Pathwise and score-function estimators}
\label{sec:rw-estimators}
When a model emits discrete symbols and the objective depends on them, the
gradient of the expected objective is classically estimated with the
score-function (REINFORCE) estimator \cite{williams1992reinforce}.  Its
group-normalised form, GRPO \cite{shao2024grpo}, is now the default for
post-training language models.  Such estimators are unbiased but their variance
grows with the number of sampled symbols \cite{mohamed2020mcgrad}, and in our
setting every sample costs one physics simulation.  The alternative is to relax
the sampling step: the Gumbel-softmax / Concrete distribution
\cite{jang2017gumbel,maddison2017concrete} combined with a straight-through
estimator \cite{bengio2013ste} makes the sampled symbol a differentiable
function of the logits, at the price of bias.

That construction has recently reached LLM post-training.  GRADE replaces policy
gradients with direct back-propagation through a Gumbel-softmax
straight-through relaxation of token sampling.  It reports gradient variance more
than $14\times$ lower than REINFORCE on sentiment-controlled generation
\cite{grade2026}.  DiffRO uses Gumbel-softmax to make a reward computed directly
on neural codec tokens differentiable, and reports word-error-rate results
rather than a variance comparison \cite{diffro2025}.  Both require the reward to
be a differentiable function of the emitted symbols, which for a text or codec
reward means a learned reward model or a differentiable decoder.  Amortized
inference goes further and makes the target a posterior rather than a reward
\cite{hu2023abs}.  In the physical-simulation literature the same bias--variance
account has been worked out in detail \cite{suh2022diffsim,onoda2026ddcg}.  That
work includes an unbiased first-order estimator for the discontinuous case
\cite{scheffler2026cfobg} and an unbiased pathwise/score-function mixture for
hybrid action spaces \cite{hpo2026}.

We keep the estimator machinery and replace the reward: our objective is a
physics solver.  The relaxed tokens must therefore be mapped to a simulable
physical system rather than to a soft embedding or a learned reward.  Two of our
results sit inside this literature and are stated here so that their status is
unambiguous.  First, the failure of a fixed $\tfrac12/\tfrac12$ mixture is not a
new phenomenon and not a claim about mixtures in general.  It is the predictable
consequence of a measured per-block relative bias greater than one, which makes
the pathwise gradient mis-scaled.  No convex combination repairs a scale
error.  Second, the affine calibration that does repair it is fitted from the
same simulations every arm already pays for, and it is an empirical device, not
an estimator with a guarantee.

\subsection{Physics gradients into learned designers}
\label{sec:rw-designers}
A parallel line feeds gradients from a differentiable simulator into the weights
of a neural designer.  GLOnet trains a generator of dielectric metasurfaces
directly on the physics gradient \cite{jiang2019glonet}.  Differentiable
programming systems make the same pattern available for mechanics and fluids
\cite{hu2020difftaichi}.  Physical design has also been posed end-to-end on top of
a \emph{learned} differentiable simulator \cite{allen2022physicaldesign}.  Learned,
that is, rather than an adjoint of a classical solver.  That distinction matters
when the guarantee one wants is fidelity to a reference engine.
The alternative route is to learn the physics instead of differentiating the
the solver, as physics-informed networks do \cite{raissi2019pinn}.  In flow control
the gradient path is now exposed by shared infrastructure rather than built per
project.  FluidGym is the first standalone, fully differentiable benchmark suite
for active flow control, written in PyTorch on the GPU-accelerated PICT solver
with PPO, SAC, DPC and TD-MPC baselines \cite{fluidgym2026}.  HydroGym is a
solver-independent reinforcement-learning platform for fluid dynamics
\cite{hydrogym2026}.  Both are platforms and benchmarks, not estimators, and
neither addresses the case in which the policy emits discrete structural
decisions.

In topology optimisation the designer's output is itself continuous.  It is a
density field, so the gradient path is closed by construction.  This holds whether
the field is produced by SIMP \cite{bendsoe1989simp} or by a meta-learned
neural-field parameterisation optimised bilevel \cite{metaneuralTO2025}.  The
most instructive recent case is \cite{topotransformer2026}, which is not a
transformer but a sensitivity-conditioned Bernoulli flow-matching generator.  It
models the topology-optimisation pipeline as a causal Markov chain.  It then uses
the data processing inequality to argue that the adjoint sensitivity
(reduced-gradient) field is an information-theoretically optimal conditioning
signal for topology prediction.  It also introduces ``pseudo-sensitivities'' for
the cases where the exact adjoint is expensive or unavailable.  This is the closest
existing instance of adjoint information reaching a learned designer, and the
distinction is exact.  In \cite{topotransformer2026} the adjoint enters as an
\emph{input field} to a trained generator.  In our setting it enters as a
\emph{gradient with respect to the designer's weights}, at training time, from a
single simulation of the design the model actually proposed.  Their result is
nonetheless an independent argument for our premise.  If the adjoint sensitivity
is the maximally informative signal about the optimal topology, then discarding it
discards precisely the information a design agent most needs.  Scalar-reward
post-training does exactly that.  Mechanically, our adjoint is an ordinary
implicit-function-theorem derivative of a fixed-point solver
\cite{giles2000adjoint,lorraine2020implicit,knoll2004jfnk}, and we claim no
part of that construction as new.

The gap we address is where this construction fails.  When the designer's output
is a sequence of categorical structural decisions, there is no continuous field
to differentiate, and the relaxation must be supplied by the design language.

\subsection{Language-model agents for physical design}
\label{sec:rw-llm}
LLM agents have been used to propose structures, write simulation code and
iterate on numerical feedback.  SGA alternates LLM-proposed discrete hypotheses
with differentiable-simulation refinement of continuous parameters, posing the
two as bilevel optimisers \cite{ma2024sga}.  Phythesis is our own prior work and
shares two authors with this paper.  It is a bi-level framework for
energy-efficient data-centre design in which an LLM level generates
simulation-ready three-dimensional layouts while a physics-guided
\emph{evolutionary} level selects assets and tunes their parameters.  It is
gradient-free by construction \cite{phythesis2025}.  The Mechanical Designer
couples a language-model agent to a numerical solver and iterates on that
solver's output \cite{jadhav2026mechdesigner}.  TextGrad propagates
natural-language feedback through a computation graph in place of numerical
gradients \cite{yuksekgonul2024textgrad}.  ChatTwin generates a data-centre
digital twin from language \cite{li2023chattwin}, and the conference version of
the present work supplies the multi-agent CFD-modelling platform whose
liquid-cooling backend we use as a production solver \cite{li2026dctwin}.

A second, newer line post-trains language models on simulator feedback.
SciDesignBench covers 520 simulator-grounded inverse-design tasks across 14
scientific domains and four settings.  The settings are single-shot design,
short-horizon feedback, long-horizon refinement and seed-design optimisation.  It
reports that the best zero-shot model reaches only $29.0$\% success on its
ten-domain shared core.  It introduces RLSF, a simulator-feedback recipe with
which a tuned 8B model raises single-turn success by $8$--$17$ percentage points
across three domains.  The authors position simulator-grounded inverse design
explicitly as a substrate for ``amortizing expensive test-time compute into model
weights'' \cite{scidesignbench2026}.  LEAPBench makes the trajectory, not the
outcome, the measured quantity.  Over 55 tasks it pairs a best-so-far
area-under-the-curve metric with a classical Bayesian-optimisation reference.  It
finds that switching from final-outcome to trajectory scoring changes the
best-model decision on $53$\% of tasks.  It also finds that LLMs do not outperform
the classical Bayesian baseline \cite{zhang2026leap}.

In all of these the simulator's influence on the model is mediated by a scalar
or by text.  Either the numbers re-enter through the prompt, in which case
nothing is amortised into the weights.  Or they enter as a reward, in which case
credit assignment falls back on a score-function estimator.  We differ in the
interface, not in the ambition: the simulator's Jacobian, not its output value,
reaches the model's parameters.  We also take \cite{zhang2026leap} as the
correct methodological warning and adopt both of its prescriptions.  We report
design quality reached per physics evaluation.  We compare against
gradient-free search rather than against another language model
(contribution~4 in \cref{sec:contributions}).

\subsection{Discrete structural design with continuous physics}
\label{sec:rw-structural}
The underlying engineering problem --- choosing a structure whose value is
revealed only after continuous physics is re-optimised --- has a mature
literature outside machine learning.  Topological derivatives quantify the
first-order effect of an infinitesimal structural perturbation
\cite{sokolowski1999topological}.  Adjoint methods have also been applied directly
to the topological design of large-scale district-heating networks under
nonlinear models \cite{blommaert2020adjoint}.  There the adjoint supplies
gradients for the continuous sizing variables while the network topology itself
is still resolved by an outer search.  That is exactly the step we make
differentiable.  Decomposition methods such as Benders'
\cite{geoffrion1972benders} and their recent learning-augmented variants
\cite{agyeman2026benders} separate a discrete master problem from a continuous
subproblem.  The hydraulic steady state we differentiate is the Global Gradient
Method \cite{todini1988gga}, solved by a damped Newton iteration of
the kind surveyed in \cite{knoll2004jfnk}.

For liquid-cooled data centres specifically, the control side is well
established: LC-Opt benchmarks reinforcement learning and agentic AI for
end-to-end liquid-cooling optimisation of a plant whose cooling topology is
given \cite{lcopt2025}.  On the design side, co-design work varies structure and
control together \cite{jadhav2026codesign,jadhav2026lifecycle}, and
system-level digital twins of liquid-cooled machines provide the surrounding
facility model \cite{brewer2024exadigit}.  In the most directly comparable of
these the structural decision is the integer partition of 25 coolant
distribution units across subloops.  It is resolved by enumerating all 611
feasible partitions against a reduced-order surrogate of a Frontier-calibrated
Modelica model, over a full year of operation \cite{jadhav2026codesign}.  These methods solve
one instance at a time, and they are the right answer when the structural space
is enumerable.  Our aim is complementary: to make the structural decision itself
differentiable.  The same physics gradients these methods already compute can then
be amortised into a model, and we say precisely when that is worth doing
(contribution~4 in \cref{sec:contributions}).

\subsection{Concurrent work: an exactly differentiable water-network solver}
\label{sec:rw-concurrent}
One result overlaps part of our solver-side construction, and we state the
overlap explicitly rather than let a reviewer find it.  Mu et al.\
\cite{mu2026gga}, posted in September 2026, make
the Global Gradient Method itself exactly differentiable.  Their forward pass
reproduces the reference engine including discrete devices, status switching and
low-flow linearisation.  Their backward pass solves the implicit adjoint by
reusing the forward pass's terminal factorisation.  One extra sparse solve
therefore returns every parameter's gradient at once, batched over scenarios on
one graphics processor.  Across 52 public, synthetic and operational networks and
8{,}140 simulation frames the largest head deviation from EPANET~2.2 is
$1.137\times10^{-13}$\,ft, and 25 networks agree exactly.  On the 905-pipe L-TOWN
benchmark one adjoint solve replaces the 906 simulations a finite-difference
roughness Jacobian costs.  Gradient calibration reaches its endpoint within a
median 595 model calls, where the strongest of five tuned metaheuristics needs
8{,}060.  This is a stronger and more thoroughly validated version of the
adjoint we attach in \cref{sec:benchB}.  We regard the construction as
established rather than as a contribution of ours: a reader who needs a
production-grade differentiable hydraulic solver should use theirs.

The two works are complementary in a way that sharpens rather than weakens our
claim.  Their decisions are roughness calibration, leak localisation and sensor
placement.  These are decisions about \emph{continuous} parameters.  Their paper
is explicit that the combinatorial part of the physics is frozen, not relaxed.
It states: ``status switching is a genuine combinatorial event: a switch changes
the effective topology or replaces an equation.  The backward pass freezes the
configuration the forward pass converged to.''  Our problem is exactly the
perturbation that would change that configuration.  The perturbations are a rack
reassigned to another coolant distribution unit, a branch added or isolated, and a
pipe class changed.
Nothing in \cite{mu2026gga} relaxes a discrete decision.  The words
\emph{language model}, \emph{Gumbel}, \emph{straight-through}, \emph{token} and
\emph{relaxation} do not occur in it.  The adjoint we attach is therefore
not a competing construction but the same construction evaluated at a point of a
design space their solver never enters.  That a verified differentiable GGA now
exists, and that it stops precisely at the combinatorial boundary, is the
clearest statement we know of the gap this paper fills.

Three further literatures bound this paper, and we name them because each owns an
ingredient.  Ground-structure topology optimisation fixes a superset of candidate
members together with a minimum size, so that every relaxed point is a physically
valid structure by construction \cite{dorn1964automatic,zegard2014grand3d,guest2004minimum};
our relaxation is that construction instantiated inside a production solver's
friction law, and what we add is not closure itself but the fidelity curve of its
floor and the invariance of node degrees, sparsity pattern and Newton path.  The
shape-derivative versus topological-derivative dichotomy
\cite{sokolowski1999topological} predicts our decision-type asymmetry: a first-order
size derivative ranks parametric edits well, while structural edits are the domain of
the topological derivative, which we do not implement.  Sec.~VII-B is therefore a
measurement of a known dichotomy in a new engineering domain, not a new dichotomy.
Neural combinatorial optimisation trains policies over discrete structures scored by a
simulator \cite{vinyals2015pointer,bello2017neural,kwon2020pomo}, and
grammar-constrained decoding enforces structured validity at generation time
\cite{hokamp2017lexically,post2018fast,willard2023guided}; both supply components we
reuse, and neither attaches a simulator adjoint to the policy's parameters.  The
deep-RL evaluation literature sets the reporting standard we adopt: seeded interval
estimates rather than single-run comparisons \cite{agarwal2021precipice,henderson2018matters},
which is why every comparison here is paired, seeded and quoted with its exact-test
floor.  The optimiser itself combines LoRA \cite{hu2021lora} with a clipped surrogate
\cite{schulman2017ppo} and truncated Gumbel sampling \cite{kool2019stochastic}.

\subsection{Positioning}
\label{sec:rw-positioning}
\Cref{tab:position} places the present work against the groups above.  Two
columns are new relative to the usual positioning table and both reflect the
corrected claims of this paper.  The column ``Relaxation is physical'' asks
whether every point of the relaxed design space is a system the solver can
simulate.  This is the property we certify, and it is not a property any estimator
paper needs.  The column ``Pathwise bias'' asks what each line of work establishes
about the bias of the pathwise estimator.  It records that the framework-level
answers are owned by
\cite{suh2022diffsim,onoda2026ddcg,scheffler2026cfobg,hpo2026} while ours is a
measurement, resolved by decision type, in one engineering domain.  The
distinguishing combination remains a \emph{discrete} structural design emitted by
a \emph{language model} whose parameters receive a \emph{solver adjoint} through a
relaxation that is itself a physical design.  Each ingredient exists; the object
that connects them does not.

% Full-width float: the table gained two columns relative to the draft's Table I.
\begin{table*}[!t]\centering
\caption{Positioning.  ``Adjoint to weights'' means the simulator's Jacobian
reaches the designer's parameters, as opposed to its output value reaching them
as a reward or as prompt text, or its sensitivity field reaching them as a
conditioning input.  ``Relaxation is physical'' means every point of the relaxed
design space is a system the solver can simulate.  ``Pathwise bias'' records what
the work establishes about the bias of a first-order/pathwise gradient estimator;
``measured by decision type'' is a measurement on one production solver, not a
bound or a correction.}
\label{tab:position}
\footnotesize
\setlength{\tabcolsep}{5pt}
\begin{tabular}{@{}lccccl@{}}\toprule
 & discrete & LLM & adjoint & relaxation & \\
 & structure & designer & to weights & is physical & pathwise bias\\
\midrule
Score-function post-training \cite{shao2024grpo,scidesignbench2026}
  & \checkmark & \checkmark & -- & n/a & none used (unbiased)\\
STE through tokens \cite{grade2026,diffro2025}
  & \checkmark$^{\dagger}$ & \checkmark & -- & -- & variance reported only\\
Amortized discrete inference \cite{hu2023abs}
  & \checkmark & \checkmark & -- & n/a & n/a (no simulator)\\
Differentiable-sim.\ policy gradients \cite{suh2022diffsim,onoda2026ddcg,scheffler2026cfobg}
  & -- & -- & \checkmark & n/a & analysed, switched, corrected\\
Unbiased mixed gradients \cite{hpo2026}
  & \checkmark & -- & \checkmark & -- & unbiased pathwise/SF mixture\\
Physics-driven designers \cite{jiang2019glonet,metaneuralTO2025}
  & -- & -- & \checkmark & n/a & n/a (continuous output)\\
Sensitivity-conditioned TO \cite{topotransformer2026}
  & -- & -- & adjoint as \emph{input} & n/a & n/a (continuous output)\\
LLM design agents \cite{ma2024sga,phythesis2025,jadhav2026mechdesigner}
  & \checkmark & \checkmark & -- & -- & none (no gradient)\\
Adjoint network design \cite{blommaert2020adjoint,mu2026gga}
  & -- & -- & \checkmark & n/a & continuous parameters only\\
Trajectory-scored design agents \cite{zhang2026leap}
  & \checkmark & \checkmark & -- & n/a & none (Bayesian reference)\\
\textbf{This work}
  & \checkmark & \checkmark & \checkmark & \checkmark & measured by decision type\\
Ground-structure TO \cite{dorn1964automatic,zegard2014grand3d,guest2004minimum}
  & \checkmark & -- & -- & \checkmark & n/a (no learning)\\
Topological derivative \cite{sokolowski1999topological}
  & \checkmark & -- & -- & n/a & structural sensitivity theory\\
Neural combinatorial optimisation \cite{vinyals2015pointer,bello2017neural,kwon2020pomo}
  & \checkmark & -- & -- & n/a & none used\\
Grammar-constrained decoding \cite{hokamp2017lexically,post2018fast,willard2023guided}
  & \checkmark & partial & -- & n/a & n/a\\
RL evaluation methodology \cite{agarwal2021precipice,henderson2018matters}
  & -- & -- & -- & n/a & reporting standard adopted here\\
\bottomrule
\multicolumn{6}{@{}l@{}}{\footnotesize$^{\dagger}$discrete text or codec tokens,
not structural design decisions.}
\end{tabular}
\end{table*}
